\documentclass[11pt]{article}

\usepackage[margin=1in]{geometry}
\usepackage{amsmath,amssymb,amsthm}
\usepackage{hyperref}

\newtheorem*{theorem}{Literature Status Theorem}

\title{Literature Status: One-Sided James Theorem}
\author{FA/Banach run \texttt{fa\_banach\_001}}
\date{2026-06-18}

\begin{document}
\maketitle

\section*{Packet Status}

This is a \texttt{literature\_already\_answered} packet, not an original
partial result from this run. It records that the arbitrary-Banach-space
version of the one-sided James compactness theorem asked for by
Cascales--Orihuela--Perez is already answered in later literature.

\section*{Original Source}

The source paper is B. Cascales, J. Orihuela, and A. Perez,
\emph{One-sided James' compactness theorem}, arXiv:1508.00496, published in
\emph{Journal of Mathematical Analysis and Applications} 445(2), 1267--1283,
2017, DOI \texttt{10.1016/j.jmaa.2016.03.080}.

The source proves the relevant two-set theorem under the additional hypothesis
that \((B_{E^\ast},w^\ast)\) is convex block compact. Its final remarks ask
whether Theorems 2 and 10 remain valid for arbitrary Banach spaces.

\section*{Supporting Literature}

Orihuela's later paper \emph{Conic James' compactness theorem},
arXiv:1610.02763, settles the Theorem 10 direction for arbitrary Banach
spaces and records that Theorem 2 remained open at that point.

The decisive later answer is Warren B. Moors, \emph{On a one-sided James'
theorem}, \emph{Journal of Mathematical Analysis and Applications} 449(1),
528--530, 2017, DOI \texttt{10.1016/j.jmaa.2016.12.019}. Crossref records
this paper with that title, author, journal, volume, issue, and page range,
and its reference list cites the Cascales--Orihuela--Perez JMAA paper
directly.

\begin{theorem}
The arbitrary-Banach-space version of the one-sided James compactness theorem
from Cascales--Orihuela--Perez is already answered in later literature.
\end{theorem}

\begin{proof}[Status verification]
The original source asks whether the theorem can be proved without the
convex-block-compactness assumption on the dual ball. Orihuela's conic James
paper handles the companion Theorem 10 direction and leaves Theorem 2 as the
remaining issue. Moors's 2017 note is a separate later paper titled
\emph{On a one-sided James' theorem}; its bibliographic metadata identifies
the exact target paper as a cited reference. The later Moors paper therefore
turns the target from a pushable partial problem into an already-known
literature answer.
\end{proof}

\section*{Preserved Partial Work}

The previous partial packet is preserved in
\texttt{partial\_reduction\_history/}. It contains useful correct reductions
and positive families, including the one-set reduction, the affine-slice
subcase, conic-dominated families using Orihuela's theorem, and frustum or
finite-dimensional cylinder examples. These should be treated as explanatory
material, not as a current open partial to push.

\section*{Access Limitation}

The anonymous ScienceDirect request available during this update returned
HTML or metadata rather than a usable local PDF for Moors's paper. The packet
therefore stores Crossref metadata for the decisive DOI instead of a local
Moors PDF. A human with institutional access can add the PDF later.

\section*{Human Review Recommendation}

Classify this target as \texttt{literature\_known} /
\texttt{literature\_already\_answered}. Do not spend another model round
trying to extend the partial result unless the goal is a self-contained
expository reproof.

\end{document}
